\documentclass[11pt,letterpaper]{article}
\usepackage{amsmath,amsthm,amssymb,geometry}
\geometry{margin=1.15in}
\usepackage[hidelinks]{hyperref}
\newcommand{\papertitle}{Autonomous proofs of the rank cases of the minimizing Sard conjecture in dimension five with GPT-6.2 Astra}
\hypersetup{pdftitle={\papertitle},pdfauthor={Renrui Chen},
  pdfkeywords={minimizing Sard conjecture, sub-Riemannian geometry, automated theorem proving, Lean 4, GPT-6.2 Astra}}
\newcounter{proofgroup}
\setcounter{proofgroup}{1}
\newtheorem{theorem}{Theorem}[proofgroup]
\newcommand{\proofsection}[1]{\stepcounter{proofgroup}\subsection{#1}}
\setlength{\emergencystretch}{1em}
\newtheorem{lemma}[theorem]{Lemma}
\newtheorem{proposition}[theorem]{Proposition}
\newcommand{\R}{\mathbb R}
\newcommand{\im}{\operatorname{im}}
\newcommand{\rank}{\operatorname{rank}}
\begin{document}
\title{\papertitle\thanks{Project repository: \url{https://github.com/Ch-RR/minimizing-sard-five}.}}
\author{Renrui Chen}
\date{November 26, 2026}
\begingroup
\renewcommand{\textasteriskcentered}{\ensuremath{\ast}}
\maketitle
\endgroup
\begin{abstract}
We describe a human-directed proof development with GPT-6.2 Astra in OpenAI's
Codex for the five-dimensional minimizing Sard conjecture. For a smooth
bracket-generating distribution of constant rank with a smooth horizontal
metric, we prove that the endpoints of singular length minimizers from each
fixed point form a set of zero smooth outer volume.
In the rank-two case the difficult endpoint differential has a two-dimensional
cokernel. A finite-rank argument for its second variation identifies a bracket
direction that can be straightened independently of the minimizing control.
Taylor expansion and a local horizontal correction then produce two-dimensional
regions of strict descent for the distance. An anisotropic differentiation
argument excludes these regions at almost every endpoint. The result concerns
attained minima and requires no completeness assumption. The three horizontal
rank cases share one formal proof. We release the Lean sources and the
verification record for 150 project modules in Lean 4.24.0.
\end{abstract}
\section{Introduction}
\paragraph{The fixed-source problem.}
Let $M$ be connected, Hausdorff, second countable and without boundary.
In a smooth local frame $(X_i)_{i=1}^m$ of the smooth subbundle $\Delta$, write
\[
 \begin{gathered}
 \dim M=5,\quad 0<m=\rank\Delta<5,\quad g\in C^\infty(\mathrm{Sym}^2\Delta^*),\quad g>0,\\
 D_j(y)=\mathrm{span}\{X_I(y):1\le |I|\le j\},\qquad
 \bigcup_{j\ge1}D_j(y)=T_yM\quad(y\in M).
 \end{gathered}
\]
For a fixed source $x$, use the unrestricted endpoint map on horizontal paths:
\[
 \begin{aligned}
 \Omega_x&=\{\gamma\in W^{1,2}([0,1],M):\gamma(0)=x,\
           \dot\gamma(t)\in\Delta_{\gamma(t)}\ \text{a.e.}\},\\
 E_x(\gamma)&=\gamma(1),\qquad L(\gamma)=\int_0^1|\dot\gamma(t)|_g\,dt,\\
 V(y)&=d(x,y)=\inf_{\beta\in E_x^{-1}(y)} L(\beta),\\
 \gamma\text{ singular}&\iff \rank D_\gamma E_x<5 .
 \end{aligned}
\]
Only the initial point is fixed when differentiating $E_x$.
\begin{theorem}\label{main}
For each $x\in M$ and each smooth positive density $\mu$ on $M$,
\[
 \mu^*\{y:\exists\gamma\in\Omega_x,\ E_x(\gamma)=y,\ 
 L(\gamma)=d(x,y),\ \rank D_\gamma E_x<5\}=0 .
\]
\end{theorem}
The quantifier over minimizing curves includes endpoints shared by regular and
singular minimizers. The exceptional set may depend on the source.
\paragraph{Related work.} Rifford
\cite{rifford} obtains the minimizing Sard property under an almost-everywhere
Goh-rank condition and, in particular, for distributions of corank at most
two in the complete setting of that paper. A rank-two distribution in dimension
five has corank three. In dimension three the strong Sard theorem of
Belotto da Silva, Figalli, Parusi\'nski and Rifford \cite{strong} covers even
nonminimizing singular curves for analytic structures.

\paragraph{Our contribution.} For a minimizing control $u$ write $A=DG(u)$ and let $\mathcal S(u)$ denote
the support condition of Lemma~\ref{support}. The proof splits as follows
outside the growth exception and the source:
\[
\begin{array}{c|l}
 \text{condition on }u & \text{consequence for }y=G(u)\\ \hline
 \mathcal S(u),\ \rank A<5 & y\in P_x(\operatorname{Crit}P_x)
                                  \quad(\text{Lemma~\ref{normal}})\\
 \neg\mathcal S(u),\ \rank A=4 & y\in N_{x,4}\quad(\text{Lemma~\ref{scalar}})\\
 \neg\mathcal S(u),\ \rank A=3 & y\in N_{x,3}\quad(\text{Proposition~\ref{two-missing}})
\end{array}
\]
The null sets are chosen before $u$. Lemma~\ref{two-point} supplies
$V<\infty$ and $V\in C^0(M)$ wherever the measure lemmas below are applied.
\section{Agent}
\paragraph{Setting.} Renrui Chen directed the investigation, specified the
fixed-source target and required the argument to retain the original
hypotheses. GPT-6.2 Astra operated through OpenAI's Codex with access to
research literature, working drafts and a local Lean installation.
Here autonomy refers to the agent generating mathematical arguments
and proof scripts within assigned tasks. The record is one interactive
investigation with human steering.

Two instructions from the author illustrate the interaction (translated from
Chinese): \textit{Do not simplify the problem or add unnecessary hypotheses}
and \textit{First formally verify the five-dimensional theorem}. The author
also supplied a research packet assembled by another AI and asked the agent
to check its suggestions critically.

\paragraph{Construction.} The agent investigated control variations that
shorten a path while moving its endpoint only at higher order. In the
rank-two branch it developed the argument identifying a line in the
two-dimensional cokernel (Proposition~\ref{line}). It then bounded the
cost of correcting the remaining endpoint error (Proposition~\ref{two-missing}).
The author requested the complete five-dimensional statement and directed the
agent to follow the earlier draft when a proposed shortcut stalled.

\paragraph{Review and formalization.} One proposed shortcut would have
replaced an anisotropic rectangle by a single ray at all scales. For
weights $(2,4)$ the available displacements can instead satisfy
$X>0$ and $X^2/4<Y<2X^2$. No fixed ray meets these boxes arbitrarily near
zero. The argument combines increments at several scales (Lemma~\ref{box}).
Codex wrote and revised the Lean declarations and proof scripts using
compiler feedback and drafted the manuscript. The author requested
explicit quantifiers, uniform estimates and the scope of the formal checks.

Huang's sum-product study~\cite{huang} reports independent trials of
a proof-generation agent. The data here describe artifacts from a single
investigation; a success rate and a per-trial token cost were not measured.

\section{Verification results}
The three rank cases share lemmas and one final formal target.
\[
\begin{array}{l|l}
\text{check} & \text{recorded outcome}\\ \hline
\text{toolchain} & \text{Lean 4.24.0, mathlib v4.24.0}\\
\text{frozen project} & 150\text{ modules},\ 2602\text{ safe declarations replayed}\\
\text{skipped recursive auxiliaries} &5;\quad 0\text{ in the target's dependency closure}\\
\text{target axiom dependencies} &\texttt{propext},\ \texttt{Classical.choice},\ \texttt{Quot.sound}
\end{array}
\]
\paragraph{What was checked.}
The target in \texttt{FiveDimensionalMinimizingSard.lean} is
\begin{center}\small\ttfamily five\_dimensional\_metric\_minimizing\_sard\end{center}
The formal path derivative is the native manifold derivative. In geometric
notation its data include
\[
 \begin{gathered}
 E_x(\gamma)=\gamma(1),\qquad
 \text{singular}\iff D_\gamma E_x\text{ is not onto},\\
 g_{ij}:y\mapsto g_y(X_i(y),X_j(y)),\qquad g_{ij}\in C^\infty,\\
 \mu(\chi^{-1}B)=\int_B\rho_\chi\,dz\quad(B\text{ Borel in the chart}),\\
 L(\gamma)=\int|\dot\gamma|_g,\qquad d(x,y)=\inf_{E_x(\beta)=y}L(\beta).
 \end{gathered}
\]
For the wider intermediate statements, the corresponding scopes are
\[
 \begin{array}{c|l|l}
  &\text{paper}&\text{formal application}\\ \hline
 \text{Lemma~\ref{box}} & d<\infty,\ k_i>1 &d=2,\ k_i\in\mathbb N,\ k_i\ge2\\
 \text{Lemma~\ref{planes}} & k,l>1&\R^3\times\R^2,\ k,l\in\mathbb N_{\ge2}\\
 \text{Lemma~\ref{scaled-box}} &d<\infty,\ k_i\in\mathbb N_{\ge2}
       &(2,2),(2,3),\ \text{scalar inversion}\\
 \text{Lemma~\ref{flags}} & m=2,\ \gamma\text{ nonconstant}
       &m=2,\ x\ne y\quad(D_3)\\
 \text{Lemma~\ref{two-point}} &1\le r<\infty &1\le r\le4
 \end{array}
\]
The fixed-prefix formula in Lemma~\ref{flags} covers nonconstant loops.
The archive includes the module and declaration references in the formula
audit, the source hashes, and the replay instructions in the ancillary README.
A supplementary module checks Proposition~\ref{two-missing} without using
the final minimizing Sard theorem. The repository gives commands to compile
the modules, replay the declarations and run the supplementary checks.

\section{Proof analysis}
We give the full argument below. The lemma and proposition numbers are
stable identifiers shared with the companion audit. Horizontal rank one
is excluded by bracket generation. For ranks two, three and four, the
final case split appears in the proof of Theorem~\ref{main}.
\proofsection{A measure lemma for strict descent}\label{measure}
In fixed linear coordinates on $\mathbb R^d$ let
$\delta_r v=(r^{k_i}v_i)_{i=1}^d$, with arbitrary fixed real
$k_i>1$. A thin anisotropic open box need not contain any straight
ray, even when all of its small dilates are available. For example, let
$d=2$, $k=(2,4)$ and $U=(1,2)^2$. A displacement $(X,Y)\in\delta_r U$
then satisfies
\[
 r^2<X<2r^2,\quad r^4<Y<2r^4,\qquad
 \tfrac14 X^2<Y<2X^2 .
\]
No fixed straight ray contains such displacements arbitrarily near zero:
on a ray with $X>0$ the ratio $Y/X$ is constant, whereas here it is
positive and tends to zero. To apply a one-dimensional differentiation
theorem we therefore construct a curve whose increments can be assembled
from the available boxes.

\begin{lemma}[One-sided infinite descent]\label{one-sided}
For a finite right-continuous function $g$ on an open interval $I$,
the set of $t$ with
$\limsup_{h\downarrow0}(g(t+h)-g(t))/h=-\infty$ has outer measure zero.
\end{lemma}
\begin{proof}
Write $D^+g(t)=\limsup_{h\downarrow0}(g(t+h)-g(t))/h$. For $j,n\ge1$ let
\[
 E_{j,n}=\left\{t:
 \begin{array}{l}
 [t,t+1/j]\subset I,\quad |g(t)|\le n,\\
 g(t+h)-g(t)\le-h\quad(0<h<1/j)
 \end{array}\right\}.
\]
Right continuity permits testing $h\in\mathbb Q$ and gives
$E_{j,n}$ Borel. Also $\{D^+g=-\infty\}\subset\bigcup_{j,n}E_{j,n}$.
Cover $I$ by countably many intervals $J$ of diameter $<1/j$. For $E=J\cap E_{j,n}$,
\[
 s,t\in E,\ s>t\ \Longrightarrow\ g(s)\le g(t)-(s-t);
 \qquad G(t)=\sup(\{-n\}\cup g(E\cap[t,\infty))).
\]
Thus $G$ is bounded, nonincreasing and $G|_E=g|_E$. At almost every
$t\in E$, $t$ is a density point of $E$ and $G'(t)\in\R$. Choose
$t_i\in E$ decreasing strictly to $t$. Then
\[
 D^+g(t)\ge\lim_{i\to\infty}\frac{g(t_i)-g(t)}{t_i-t}
       =\lim_{i\to\infty}\frac{G(t_i)-G(t)}{t_i-t}=G'(t)>-\infty .
\]
Consequently $|E\cap\{D^+g=-\infty\}|^*=0$ for every $j,n,J$.
\end{proof}

\begin{lemma}[Anisotropic descent boxes]\label{box}
For continuous $V$ on an open set, the following base points $z$
form a set of Lebesgue outer measure zero. There are a nonempty
open box $U$, $c>0$ and $r_0>0$ such that for $0<r<r_0$,
\[
          V(z+\delta_r b)\le V(z)-cr \qquad(b\in U).
\]
Countably many adapted coordinate systems, translations, and
weight lists may be used. This includes arbitrarily narrow boxes.
\end{lemma}
\begin{proof}
By shrinking use rational $c,r_0$ and
$U=a+(-e,e)^d$ with rational $a,e$. Work on buffered subcharts,
extending $V$ continuously locally via cutoffs when needed.
With $\mathcal R=(0,r_0)\cap\mathbb Q$ and $\mathcal B=U\cap\mathbb Q^d$,
continuity gives the Borel test
\[
 E=\bigcap_{r\in\mathcal R,\ b\in\mathcal B}
       \{z: V(z+\delta_r b)-V(z)\le-cr\}.
\]
For $\kappa=\sum_i k_i$ define
\[
 d_k(z,w)=\max_i |z_i-w_i|^{1/k_i},\quad B_r(z)=z+\delta_r(-1,1)^d,\quad
 |B_{2r}(z)|=2^\kappa|B_r(z)|.
\]
Each coordinate satisfies the triangle inequality since $1/k_i<1$.
Lebesgue differentiation for this doubling ball basis gives, for almost every $z\in E$,
\[
 \lim_{r\downarrow0}\frac{|B_r(z)\setminus E|}{2^d r^\kappa}=0 .
\]
For fixed $\epsilon,C>0$, a center $w\in B_{Cr}(z)$ satisfies
$B_{\epsilon r}(w)\subset B_{(C+\epsilon)r}(z)$. At a density point
\begin{equation}\label{density-hit}
 |B_{(C+\epsilon)r}(z)\setminus E|=o(r^\kappa)
       < 2^d\epsilon^\kappa r^\kappa=|B_{\epsilon r}(w)|.
\end{equation}
The threshold on $r$ is independent of $w$, so the small box meets $E$.

Choose $1/\min_i k_i<\alpha<1$. In $(0,1)$ place one atom of
vector mass $\delta_{r_j}a$, $r_j=2^{-\alpha j}$, at each dyadic
rational of exact denominator $2^j$, $j\ge1$. If $N_j(t)$ counts those
rationals in $(0,t]$, the cumulative sum $D$ is given coordinatewise by
\begin{equation}\label{clock-series}
 D_i(t)=a_i\sum_{j\ge1}N_j(t)2^{-\alpha j k_i},\qquad
 |a_i|\sum_{j\ge1}2^{j-1}2^{-\alpha j k_i}<\infty .
\end{equation}
The bound is uniform in $t$, so $D$ is right-continuous. For
$0<t<t+h<1$ let $J$ be the least order of an
atom in $(t,t+h]$. There is exactly one atom at that order, and
$h\le 2^{1-J}$ since the interval contains no coarser grid point.
Put $n_m=N_{J+m}(t+h)-N_{J+m}(t)$. The spacing of the atoms gives
\[
 n_0=1,\qquad n_m\le1+h2^{J+m-1}\le 1+2^m\le 3\cdot 2^m .
\]
Fix $L$ large enough that for every $i$ the discarded tail satisfies
\[
 \begin{aligned}
 |\mathrm{tail}_i|
 &\le |a_i|\sum_{m>L}n_m r_{J+m}^{k_i}
 \le 3|a_i|r_J^{k_i}\sum_{m>L}2^{-(\alpha k_i-1)m}
 <e r_J^{k_i}/4 .
 \end{aligned}
\]
Reorder the retained scales $s_1,\ldots,s_N$ with the largest last. Then
\[
 N\le3\sum_{m=0}^L2^m,\qquad 2^{-\alpha L}r_J\le s_\ell\le s_N=r_J,\qquad
 \left|\sum_{\ell\le j}a_i s_\ell^{k_i}\right|\le N|a_i|r_J^{k_i}.
\]

At a density point $z\in E$, define $q_0=z$,
$q_j=z+\sum_{\ell=1}^j\delta_{s_\ell}a$, and choose a fixed $\epsilon>0$ with
\[
 2\epsilon^{k_i}<\tfrac14 e\,2^{-\alpha L k_i}\quad(1\le i\le d).
\]
The preceding partial-sum bound gives $q_j\in B_{Cr_J}(z)$ for a fixed $C$.
By \eqref{density-hit}, for sufficiently small $r_J$ choose
\[
 \begin{aligned}
 p_0&=q_0,\quad p_j\in E\cap B_{\epsilon r_J}(q_j)\quad(0<j<N),\\
 p_N&=z+D(t+h)-D(t),\quad \eta_j=p_j-q_j,\\
 \eta_0&=0,\quad |\eta_{j,i}|<(\epsilon r_J)^{k_i}\ (0<j<N),\quad
 \eta_N=\mathrm{tail}.
 \end{aligned}
\]
All $p_j$ for $j<N$ belong to $E$ by \eqref{density-hit}. For $j<N$,
\[
 |(p_j-p_{j-1})_i-s_j^{k_i}a_i|
 \le 2(\epsilon r_J)^{k_i}< e s_j^{k_i}.
\]
For $j=N$ the corresponding bound is
$e r_J^{k_i}/4+(\epsilon r_J)^{k_i}<e s_N^{k_i}$.
Thus each increment lies in its own admissible box. The final point need
not belong to $E$: only starting points enter the descent hypothesis.
Summing the $N$ inequalities gives
\begin{equation}\label{clock-sum}
 V(z+D(t+h)-D(t))-V(z)
 =\sum_{j=1}^N (V(p_j)-V(p_{j-1}))
 \le -c\sum_{j=1}^N s_j\le -c r_J .
\end{equation}
For each fixed $t$, $J\to\infty$ as $h\downarrow0$ and
$r_J/h\ge 2^{(1-\alpha)J-1}\to\infty$.

For every translate $z_0+D(t)$ the function $V(z_0+D(t))$ is
right-continuous and finite on a buffered chart (or use the local
extension). Wherever $z=z_0+D(t)$ is a density point of $E$ its
upper right Dini derivative is $-\infty$. Lemma~\ref{one-sided}
excludes this for almost every $t$. For the more general
Denjoy--Young--Saks alternatives see \cite[Remark 2.10]{rifford}.
Let $E^{(1)}$ denote these density points of $E$. The incidence function
$(z_0,t)\mapsto\mathbf1_{E^{(1)}}(z_0+D(t))$ is Borel. Its time integral
vanishes for each $z_0$ by the preceding derivative estimate. Tonelli's
theorem and translation invariance now give
\begin{equation}\label{incidence}
 \begin{aligned}
 0&=\int_{\R^d}\int_0^1 \mathbf1_{E^{(1)}}(z_0+D(t))\,dt\,dz_0\\
  &=\int_0^1 |E^{(1)}-D(t)|\,dt=|E^{(1)}|.
 \end{aligned}
\end{equation}
The differentiation theorem gives $|E\setminus E^{(1)}|=0$.
Taking the countable union over the rational choices finishes the proof.
\end{proof}

\begin{lemma}[A common exceptional set for fixed planes]\label{planes}
Fix a continuous $V:M\to\R$ and a countable collection of smooth charts.
In each chart prescribe countably many ordered pairs $e,f$ of distinct
coordinate axes and pairs of weights $k,l>1$. There is a Borel set $N$
of zero smooth volume containing every point $y$ at which one of these
choices gives an open rectangle $U$ and constants $c,r_0>0$ satisfying
\[
 V(y+r^k a e+r^l b f)\le V(y)-cr
 \quad((a,b)\in U,\ 0<r<r_0).
\]
Addition is in the chosen chart. The rectangle and constants may depend
on the point and on any curve witnessing the inequality.
\end{lemma}
\begin{proof}
Fix a chart $\chi$, two axes and weights $(k,l)$. Cover its target by
countably many product boxes $B$ with compact closure inside it. With a
smooth cutoff equal to one near $\overline B$, extend the coordinate
value to $\widetilde V\in C^0(\R^5)$. At a witness point shrink $U,r_0$
to stay in $B$. For fixed bounded rational $U$ and positive rational $c,r_0$,
\[
 E_\lambda=B\cap
 \bigcap_{\substack{(a,b)\in U\cap\mathbb Q^2\\r\in(0,r_0)\cap\mathbb Q}}
 \{z:\widetilde V(z+r^kae+r^lbf)-\widetilde V(z)\le-cr\}.
\]
Continuity makes this Borel test equivalent to testing all real parameters.
In the following display $V$ denotes the extended value. In coordinates $(w,z)\in\R^3\times\R^2$,
\[
 \begin{aligned}
 (E_\lambda)_w&\subset\{z:V(w,\cdot)\text{ satisfies Lemma~\ref{box} at }z\},\\
 |E_\lambda|_5&=\int_{\R^3}|(E_\lambda)_w|_2\,dw=0,\\
 N&=\bigcup_{\lambda}\chi_\lambda^{-1}(E_\lambda),\quad
 \mu(N)\le\sum_{\lambda}\int_{E_\lambda}\rho_\lambda(\zeta)\,d\zeta=0 ,
 \end{aligned}
\]
where $\rho_\lambda$ is the coordinate density. The union ranges over a
countable set; every admissible rectangle contains a rational test rectangle.
\end{proof}

\proofsection{The endpoint in control coordinates}\label{controls}
Fix a finite subdivision and an orthonormal frame $X_i^j$ on each buffered
piece. Rescale each piece to $[0,1]$. In terminal coordinates,
\[
 \begin{aligned}
 \mathcal U&=\prod_{j=1}^N L^2([0,1],\R^m),&
 \dot q_j&=\sum_i u_{j,i}X_i^j(q_j),\\
 q_1(0)&=x,\quad q_{j+1}(0)=q_j(1),&
 G(u)&=q_N(1),\\
 J(u)&=\sum_j\int_0^1|u_j|,&
 |J(u)-J(v)|&\le\sqrt N\,\|u-v\|_{\mathcal U}.
 \end{aligned}
\]
At a minimizer $V(G(u))=J(u)$, while $V(G(v))\le J(v)$ throughout the chart.
\begin{lemma}[Control charts and the endpoint derivative]\label{control-coordinates}
At every $\gamma\in\Omega_x$, the split controls of a fixed buffered
frame subdivision give smooth local coordinates. In these coordinates
$E_x$ is represented by $G$, its derivative has the rank of $DG(u)$,
and its length is $J(u)$.
\end{lemma}
\begin{proof}
On $[a,a+d]$ the rescaling and its inverse are
\[
 u^\sharp(s)=d\,u(a+ds),\qquad u(t)=d^{-1}u^\sharp((t-a)/d),
 \qquad \|u^\sharp\|_2^2=d\int_a^{a+d}|u(t)|^2\,dt .
\]
On a smaller chart tube apply a smooth cutoff to the fields and set
\[
 \mathcal T(q,z,u)(t)=z+\sum_i\int_0^t u_i(s)X_i(q(s))\,ds .
\]
Subdivide until $\kappa:=\sum_i\|DX_i\|_\infty\|u_i\|_1<1$, also for
controls near the reference control. Then
\[
 \begin{aligned}
 \|\mathcal T(q,z,u)-\mathcal T(\tilde q,z,u)\|_\infty
       &\le\kappa\|q-\tilde q\|_\infty,\\
 D_q^j\mathcal T[h_1,\ldots,h_j](t)
       &=\sum_i\int_0^t u_i D^jX_i(q)[h_1,\ldots,h_j]\,ds
                  \quad(j\ge1),\\
 (I-D_q\mathcal T)^{-1}&=\sum_{n\ge0}(D_q\mathcal T)^n .
 \end{aligned}
\]
The mixed derivatives with one control variation replace $u_i$ by that
variation; higher control derivatives vanish. The implicit function theorem
at $q=\mathcal T(q,z,u)$ gives smooth $(z,u)\mapsto q$ in $C^0$.
Writing $X(q)$ for the matrix of frame columns in a chart gives
\[
 \dot q=X(q)u\in L^2,\quad
 X^\dagger=(X^\top X)^{-1}X^\top,\quad u=X^\dagger(q)\dot q .
\]
On common subdivisions a change of coordinates $\phi$ and frame $\widetilde X$
therefore acts by
\[
 \widetilde u=\widetilde X^\dagger(\phi(q))D\phi(q)X(q)u .
\]
Composition on buffered tubes is smooth in $C^0$ and multiplication
$C^0\times L^2\to L^2$ is continuous bilinear. The reverse formula and
uniqueness of the ODE give mutually inverse smooth transitions.

Let $\eta_h$ denote the state variation along $u+\epsilon h$.
Set $B(t)=\sum_i u_i(t)DX_i(q(t))$ in each chart. Differentiating the
integral equation gives
\begin{equation}\label{linearstate}
 \dot\eta_h(t)=B(t)\eta_h(t)+\sum_i h_i(t)X_i(q(t)),\qquad \eta_h(0)=0 .
\end{equation}
Let $U(t,s)$ be the evolution operator of $\dot\eta=B\eta$, including the
derivatives of the chart transitions at intermediate nodes. Taking the
terminal coordinate derivative as part of $U(T,s)$ gives, on the concatenated
time interval $[0,T]$,
\begin{equation}\label{firstvariation}
 DG(u)h=\int_0^T U(T,s)\sum_i h_i(s)X_i(q(s))\,ds .
\end{equation}
Finally, orthonormality gives $|\dot q(t)|_g=|u(t)|$ almost everywhere on
each piece. Integrating and summing identifies the costs.
\end{proof}

\begin{lemma}[Separation and a uniform length decrease]\label{support}
Put $A=DG(u)$, $W=\im A$, and $K=\ker A$. There exists a terminal covector
$p$ satisfying $J(v)\ge J(u)+pA(v-u)$ for every $v$ if and only if $u$
minimizes $J$ on $u+K$. In its absence there is a relatively open nonempty
set of $h\in K$ with $J(u+h)<J(u)$. Near each such $h$, any smooth
kernel slice $S$ with $S(0)=u$ and $DS(0)$ the inclusion satisfies
$J(S(th'))\le J(u)-ct$ for all $h'$ near $h$ and all small $t>0$,
for some $c>0$.
\end{lemma}
\begin{proof}
If $J(u)=\inf_{u+K}J$, separate $\{(v,s):s>J(v)\}$ from
$(u+K)\times\{J(u)\}$. For some $(\ell,\lambda)\ne(0,0)$,
\[
 \ell(v)+\lambda s\ge\ell(u+k)+\lambda J(u)
 \quad(s>J(v),\ k\in K).
\]
Let $s\to+\infty$ and vary $k$ along lines: $\lambda\ge0$, $\ell|_K=0$.
If $\lambda=0$, varying $v$ forces $\ell=0$. Thus $\lambda>0$, and
\[
 p(Av)=-\ell(v)/\lambda,\qquad
 s\downarrow J(v)\ \Longrightarrow\ J(v)-J(u)\ge pA(v-u).
\]
Extend $p$ from $W$ to the target. Conversely $v\in u+K$ gives $pA(v-u)=0$.
Strict inequality persists
in a neighborhood by continuity. Choose $\delta>0$ so that
$J(u+h')\le J(u)-2\delta$ there. For $0<t<1$ convexity gives
\begin{equation}\label{slice-saving}
 \begin{aligned}
 J(u+th')&\le(1-t)J(u)+tJ(u+h')\le J(u)-2\delta t,\\
 J(S(th'))&\le J(u+th')+\operatorname{Lip}(J)\|S(th')-u-th'\|
             \le J(u)-\delta t .
 \end{aligned}
\end{equation}
Indeed, for $h'$ in a smaller ball with $\|h'\|\le R$,
\[
 \sup_{h'}t^{-1}\|S(th')-u-th'\|
 \le R\sup_{0<\|k\|\le tR}
       \frac{\|S(k)-S(0)-DS(0)k\|}{\|k\|}
 \longrightarrow0 .
\]
\end{proof}

\begin{lemma}[Kernel slice and change of target chart]\label{kernel-slice}
For any complement $C$ to $W$ in terminal coordinates there is a smooth
map $S$ on a neighborhood of $0$ in $K$, with $S(0)=u$, $DS(0)$ the
inclusion and $G(S(k))-y\in C$. Its second differential modulo $W$ is
the intrinsic Hessian $H$ on $K$. The quadratic Taylor coefficient in
direction $h$ is $\tfrac12 H(h,h)$. Under a change of terminal chart
both this tensor and the length support condition transform by the
terminal linear isomorphism.
\end{lemma}
\begin{proof}
Choose a finite-dimensional control subspace $L$ on which $A$ is an
isomorphism onto $W$, and let $\Pi$ denote projection onto $W$ parallel
to $C$. Solve $\Pi(G(u+k+\ell)-y)=0$ for $\ell\in L$.
The partial derivative in $\ell$ is $A|_L$ and the partial derivative
in $k$ vanishes, so $\ell(0)=0$, $D\ell(0)=0$. Put $S(k)=u+k+\ell(k)$.
For the quotient map $\pi:T_yM\to T_yM/W$ and $h,k\in K$ the chain rule reads
\begin{equation}\label{slice-hessian}
 \begin{aligned}
 D^2(G\circ S)(0)[h,k]&=D^2G(u)[h,k]+A D^2S(0)[h,k],\\
 \pi D^2(G\circ S)(0)[h,k]&=\pi D^2G(u)[h,k].
 \end{aligned}
\end{equation}
For $B=D\phi(y)$ the chart-change formulas are
\[
 \begin{aligned}
 \ker(BA)&=K,\quad \im(BA)=BW,\quad \pi A=0,\\
 D^2(\phi\circ G)[h,k]&=BD^2G[h,k]+D^2\phi[Ah,Ak]=BD^2G[h,k]
                       \quad(h,k\in K),\\
 [BD^2G[h,k]]_{BW}&=\bar B[D^2G[h,k]]_W,\qquad (pB^{-1})(BA)=pA ,
 \end{aligned}
\]
where $\bar B$ is the isomorphism on the quotients.
\end{proof}

\begin{lemma}[Uniform inversion after scaling]\label{scaled-box}
Let $F$ be smooth near $0$ in $K$, with values in $\R^d$. For integers
$k_i\ge2$ suppose all derivatives of $F_i$ of orders $0,\ldots,k_i-1$
vanish at zero. Write $P_i(h)=D^{k_i}F_i(0)[h^{k_i}]/k_i!$, where
$[h^{k_i}]$ denotes $k_i$ copies of $h$. If $DP(h)$ is onto at a
direction $h$ in an open set $O\subset K$, then for each neighborhood
$O'\subset O$ of $h$ there are an open box $B$ about $P(h)$ and $t_0>0$
such that, for every $0<t<t_0$ and every $a\in B$, some $h'\in O'$
satisfies $F_i(th')=t^{k_i}a_i$ for all $i$.
\end{lemma}
\begin{proof}
Choose $d$ directions defining a linear map $L:\R^d\to K$ with
$DP(h)L$ invertible. Write $h'=h+Lz$. Taylor's integral formula extends
$t^{-k_i}F_i(t h')$ smoothly across $t=0$ on this finite-dimensional family:
\[
 \frac1{(k_i-1)!}\int_0^1(1-s)^{k_i-1}
       D^{k_i}F_i(sth')[h'^{\,k_i}]\,ds .
\]
At $t=0$ this is $P_i(h')$. Denote the extended vector by $H(t,z)$.
For $\Psi(t,z)=(t,H(t,z))$,
\[
 D\Psi(0,0)=
 \begin{pmatrix}1&0\\ \partial_t H(0,0)&DP(h)L\end{pmatrix}.
\]
The inverse theorem gives $(-t_0,t_0)\times B\subset\operatorname{dom}\Psi^{-1}$
near $(0,P(h))$. Shrink this product so that for $(t,z)=\Psi^{-1}(t,a)$
we have $h+Lz\in O'$. Then $F_i(t(h+Lz))=t^{k_i}a_i$ for every $i$.
\end{proof}

\begin{lemma}[Exclusion of full polars and unsupported corank one]\label{scalar}
Outside a common null set of endpoints every singular minimizing control with
$\rank A=4$ has a supporting covector. If $\rank A=3$, then outside a common
null set any unsupported minimizing control satisfies
\[
 \rank\big(k\mapsto [D^2G(u)(h,k)]\big)<2\quad(h\in K).
\]
\end{lemma}
\begin{proof}
In rank four choose a complementary coordinate axis $e$. On a strictly
shorter slice ray, write $G(S(th))=y+f(t)e$, $|f(t)|\le Ct^2$.
The saving $ct$ precludes $f(t)=0$ for $0<t\le T$. Let
$\sigma=\operatorname{sgn} f(T)$. By continuity, for $0<r<|f(T)|$,
\[
 \begin{aligned}
 &\exists\,t\in(0,T]:\quad f(t)=\sigma r,\quad r\le Ct^2,\\
 &\frac{V(y+\sigma r e)-V(y)}r\le-\frac{ct}{r}
                         \le-\frac c{\sqrt{Cr}}\longrightarrow-\infty .
 \end{aligned}
\]
For each signed axis the Dini condition is Borel by rational increments.
Lemma~\ref{one-sided} on each parallel line gives
$\int_{\R^4}|\{s:D^+V(w,s)=-\infty\}|\,dw=0$ in signed axis coordinates.
Take the countable union over charts and axes.

For $\rank A=3$ choose complementary axes. Put $H=\pi D^2G(u)|_{K\times K}$.
An onto polar at $h_*$ gives
\[
 b_1,b_2\in K,\qquad M(v)=\det[H(v,b_1),H(v,b_2)],\qquad M(h_*)\ne0 .
\]
For $h_0$ in the open shorter set, $M(h_0+\lambda h_*)$ has leading
coefficient $M(h_*)$ in degree two. Choose a sufficiently small $\lambda$
outside its zero set, so $h=h_0+\lambda h_*$ is shorter. Now
\[
 Q(v)=\tfrac12 H(v,v),\quad DQ(h)b=H(h,b),\quad
 \det[DQ(h)b_1,DQ(h)b_2]=M(h)\ne0 .
\]
In slice coordinates $F=G\circ S-y$ with complementary plane identified
with the quotient. Lemmas~\ref{scaled-box} and~\ref{support} give
\[
 \begin{gathered}
 \exists\, B\text{ open box},\ c>0,\ t_0>0\quad
 \forall a\in B,\ 0<t<t_0\quad\exists h'\text{ near }h:\\
 F(th')=t^2a,\qquad J(S(th'))\le V(y)-ct ,
 \end{gathered}
\]
where $B$ is a nonempty open box in the plane. This is the $(2,2)$
test of Lemma~\ref{planes}. Take the finite union over complementary axis
choices in each chart and the countable union over charts. The continuity
of $V$ used here follows from Lemma~\ref{two-point}.
\end{proof}
\proofsection{Terminal brackets and the intrinsic Hessian}\label{brackets}
Let $D_j$ denote the pointwise span of brackets of length at most
$j$ in any local smooth frame.
\begin{lemma}[Growth outside a common null set]\label{regular-growth}
There is a Borel null set outside which, for every $j\ge1$ with
$D_j(y)\ne T_yM$, one has $\dim D_{j+1}(y)>\dim D_j(y)$.
Consequently $m=2$ implies $\dim D_2=3$, $\dim D_3\ge4$, $D_4=TM$
there; and $m=3$ implies $\dim D_2\ge4$, $D_3=TM$ there.
\end{lemma}
\begin{proof}
Fix $j$, $r<5$ and a subchart with independent words
$B_1,\ldots,B_r$ of orders $\le j$. Complete them to a smooth frame
$(B_a)_{a=1}^5$. On this subchart let
\[
 E=\{z:\dim D_j(z)=\dim D_{j+1}(z)=r\},\qquad
 X_I=\sum_{a=1}^5 b_{I,a}B_a .
\]
Then $b_{I,a}|_E=0$ for $a>r$ and $|I|\le j+1$. For a smooth $b$ vanishing
on $E$, a nonzero $Db(y)$ at a density point $y\in E$ would give a cone of
directions $\omega$ of positive solid angle and some $\epsilon>0$ with
\[
 |b(y+s\omega)|\ge \epsilon s\quad(0<s<s_0).
\]
Such a cone contradicts density one. Applying this argument to each partial
derivative on the preceding full-measure subset gives, by countable intersection,
\[
 \partial^\alpha b_{I,a}(y)=0\quad
 (|\alpha|\ge0,\ |I|\le j+1,\ a>r),\qquad\text{a.e. }y\in E .
\]
At one such $y$, let $\mathcal F_y$ be the ideal of flat smooth germs and
$\mathcal V$ the module of all vector field germs. Set
\[
 \mathcal L=\sum_{a\le r} C^\infty_y B_a+\mathcal F_y\mathcal V.
\]
The preceding identities imply $X_i,[X_i,B_a]\in\mathcal L$ for $a\le r$.
The product rule gives
\[
 [X_i,fB_a]=(X_i f)B_a+f[X_i,B_a],\quad
 X_i\mathcal F_y\subseteq\mathcal F_y,\quad
 [X_i,\mathcal L]\subseteq\mathcal L .
\]
Induction gives all brackets nested with a generator; Jacobi and bilinearity
express every bracket word as a linear combination of such brackets. Hence
\[
 T_yM=\bigcup_nD_n(y)\subseteq\mathcal L(y)
       =\mathrm{span}\{B_a(y):a\le r\},\qquad 5\le r ,
\]
a contradiction. The rank sets are Borel (smooth minors); countably many
$j,r$, frame charts and choices of words cover the exceptions.
Thus for rank two in dimension five, outside a fixed null set,
\[
       \dim D_2=3,\qquad \dim D_3\ge4,\qquad D_4=TM.       \tag{G}\label{growth}
\]
The equality for $D_2$ uses that there is only one first bracket.
At $m=3$ strict growth from dimension three gives $\dim D_2\ge4$
and $D_3=TM$. Enlarge the countable union of exceptions to a Borel null set.
\end{proof}

For two frames $X$ and $\widetilde X=UX$, $U\in GL_m(C^\infty)$, use
\[
 [aX_i,bX_j]=ab[X_i,X_j]+a(X_ib)X_j-b(X_ja)X_i .
\]
Induction gives $D_j^{\widetilde X}\subseteq D_j^X$; using $U^{-1}$ gives equality.
\begin{lemma}[Terminal vanishing]\label{causal}
Let $T:L^2(0,1)\to L^2(0,1)$ be a finite-rank integral operator
with bounded kernel $\mathbf1_{s<t}K(t,s)$, continuous on $s<t$.
There exists $\tau<1$ such that $K(t,s)=0$ whenever $\tau<s<t<1$.
\end{lemma}
\begin{proof}
Write $\mathcal H_a$ for the inputs zero almost everywhere before $a$,
and put $R_a=T(\mathcal H_a)$. Nesting the domains gives
$R_b\subset R_a$ for $a<b$, and causality gives $R_a\subset\mathcal H_a$.
Choose a rational $\tau\in(0,1)$ minimizing $\dim R_\tau$ among
rational cuts. Finite dimensionality gives $R_b=R_\tau$ for every
rational $b\in(\tau,1)$, and hence
\begin{equation}\label{tail-rank}
 R_\tau\subset\bigcap_{b\in\mathbb Q\cap(\tau,1)}\mathcal H_b=\{0\}.
\end{equation}
For $\tau<s<t<1$, take intervals $I_\epsilon\ni s$, $J_\epsilon\ni t$
of diameter $\le\epsilon$ with $\tau<I_\epsilon<J_\epsilon<1$. Since
$T{\bf1}_{I_\epsilon}=0$,
\[
 0=\frac{\langle{\bf1}_{J_\epsilon},T{\bf1}_{I_\epsilon}\rangle}
          {|I_\epsilon||J_\epsilon|}
 =\frac{\int_{J_\epsilon}\int_{I_\epsilon}K(v,w)\,dw\,dv}
        {|I_\epsilon||J_\epsilon|}\longrightarrow K(t,s).
\]
\end{proof}

\begin{lemma}[Terminal bracket constraints]\label{flags}
The image $W$ contains $\Delta_y$ for every path. Suppose in addition
that $m=2$ and the path is nonconstant. Then $D_2(y)\subset W$, and
every covector $p$ annihilating $W$ for which $pD^2G(u)$ has finite
rank on $K\times K$ annihilates $D_3(y)$.
\end{lemma}
\begin{proof}
For $p\in W^\perp$ (the annihilator), put $p(t)=pU(T,t)$. By
\eqref{firstvariation}, on each piece,
\[
 0=pAh=\sum_i\int p(t)X_i(q(t))h_i(t)\,dt\quad(\forall h)
 \ \Longrightarrow\ p(t)X_i(q(t))=0 .
\]
The coefficients are continuous, so the equality holds also at the
terminal time. Thus $W^\perp\subseteq\Delta_y^\perp$ and $\Delta_y\subseteq W$.
For a nonconstant loop choose $s<1$ with $\gamma(s)\ne y$.
In a subdivision
containing $s$, let $I_s$ insert a terminal variation by zero on the
prefix. In the same terminal chart the endpoint germ $G_s$ of this
suffix at its reference input $u_s$ satisfies
\[
 G_s(u_s+v)=G(u+I_sv),\quad
 DG_s=DG\,I_s,\quad D^2G_s[h,k]=D^2G[I_sh,I_sk].
\]
In particular, with $K_s=\ker DG_s$ and $B_p=(pD^2G)|_{K\times K}$,
\[
 \begin{gathered}
 \im DG_s\subset W,\qquad I_sK_s\subset K,\\
 (pD^2G_s)|_{K_s\times K_s}=(I_s|_{K_s})^*B_p,\qquad
 \rank (I_s|_{K_s})^*B_p\le\rank B_p .
 \end{gathered}
\]
Discard stationary final pieces. On the last moving piece in unit time put
\[
 a=\min\{s\in[0,1]:\int_s^1|u(t)|\,dt=0\}.
\]
Continuity of the integral gives $a>0$, $q(a)=y$, $u=0$ a.e.\ on $[a,1]$.
Extend variations by zero there: the varied endpoint at $a$ equals that at
$1$, so the same kernel and rank restriction applies. Rescaling $[0,a]$ gives
\[
 \widehat u(t)=a\,u(at),\qquad
 \int_\sigma^1|\widehat u(t)|\,dt=\int_{a\sigma}^a|u(t)|\,dt>0
                  \quad(0\le\sigma<1).
\]

On a short final interval write $q$ for its initial state and $\Phi_t$
for the reference flow from that time. For $X_u=\sum_i u_iX_i$ put
\[
 A_t=D\Phi_t,\qquad Y_i(t)=\Phi_t^*X_i,\qquad B(t)=\Phi_t^*[X_1,X_2].
\]
On compact buffered tubes the differentiated flow equations give
$\|\partial_tD^j\Phi_t\|_\infty\le C_j\sum_i |u_i(t)|$ a.e.:
differentiate $X_u\circ\Phi_t$ in space and use bounded lower spatial
derivatives inductively and Gronwall for the highest one. Hence
\[
 \begin{aligned}
 A_t'&=(DX_u\circ\Phi_t)A_t,\qquad (A_t^{-1})'=-A_t^{-1}(DX_u\circ\Phi_t),\\
 \partial_t(\Phi_t^*Z)&=\partial_t\{A_t^{-1}(Z\circ\Phi_t)\}\\
 &=A_t^{-1}\big((DZ\,X_u-DX_u\,Z)\circ\Phi_t\big)=\Phi_t^*[X_u,Z]\quad\text{a.e.}
 \end{aligned}
\]
The identical calculation differentiated once in the initial
state is justified by the bounds above on the differentiated integral
equations. Taking $Z=X_1,X_2$ therefore gives the first spatial jets
\[
  (j^1Y_1)'=-u_2j^1 B,\qquad (j^1Y_2)'=u_1 j^1B .
\]
Here $p_0$ is the transported endpoint row at $q$. Since $p_0Y_i(t,q)=0$,
\[
 u_i(t)p_0B(t,q)=0\ \text{a.e.}\ (i=1,2).
\]
If $p_0B(1,q)\ne0$, then $|p_0B(t,q)|\ge|p_0B(1,q)|/2$ on a final
interval, contradicting $\int_\sigma^1|u|>0$. Hence
$p[X_1,X_2](y)=0$ for every $p\in W^\perp$, proving $D_2(y)\subseteq W$.

For the third layer pass to the moving equation
$\dot z=\sum v_i(t)Y_i(t,z)$ based at $q$, with zero reference perturbation.
With $z(0)=q$, the change of variables is
\[
 \begin{aligned}
 \frac d{dt}\Phi_t(z(t))
 &=X_u(t,\Phi_t(z(t)))+D\Phi_t(z(t))\sum_i v_i(t)Y_i(t,z(t))\\
 &=\sum_i(u_i(t)+v_i(t))X_i(\Phi_t(z(t))) .
 \end{aligned}
\]
For the a.e.\ chain rule on a compact buffered tube, in the uniform norm,
\[
 \Phi_t-\Phi_s=\int_s^t X_{u(\tau)}\circ\Phi_\tau\,d\tau,\qquad
 \|X_{u(t)}\circ\Phi_t\|_\infty\le C\sum_i|u_i(t)|.
\]
This Bochner integral gives time differentiation on a common full-measure
set. Boundedness and continuity of $D\Phi_t$ then give the chain rule
along the absolutely continuous $z(t)$. Let $F$ denote the moving endpoint,
$\psi$ the final reference flow in terminal coordinates, and $I$ insertion of
the terminal perturbation (zero elsewhere, with the inverse time rescaling).
Uniqueness of solutions gives $G(u+Iv)=\psi(F(v))$ near $v=0$. Put
$p_0=pD\psi(F(0))$. Differentiating once gives
$I(\ker DF(0))\subset K$; twice, for $h,k\in\ker DF(0)$, gives
\[
 p_0D^2F(0)[h,k]=pD^2G(u)[Ih,Ik],
\]
since $D^2\psi[DF(0)h,DF(0)k]=0$. Pullback of a finite-rank form
has finite rank. Thus the moving scalar Hessian has finite rank on the moving kernel.
Put $K_0=\ker DF(0)$, choose a complement of dimension
$r_0=\rank DF(0)$, and denote the projection to $K_0$ by $P$, $Q=I_{\rm id}-P$,
with $I_{\rm id}$ the identity. For the full scalar Hessian $B$,
\[
 B=P^*BP+Q^*BP+BQ,\qquad
 \rank B\le\rank(B|_{K_0\times K_0})+2r_0<\infty .
\]

The moving integral equation and its first variation at $v=0$ are
\[
 \begin{aligned}
 z_v(t)&=q+\int_0^t\sum_i v_i(s)Y_i(s,z_v(s))\,ds,\\
 H_h(t)&=\int_0^t\sum_j h_j(s)Y_j(s,q)\,ds .
 \end{aligned}
\]
In the mixed derivative at zero a variation must hit the explicit control
factor for a term to survive. Hitting it with $k$ leaves $D Y_i\,H_h$;
hitting it with $h$ leaves $D Y_i\,H_k$. Pairing the time-one value with
$p_0$ gives the scalar Hessian
\[
 \int_0^1\sum_i p_0DY_i(t,q)
 \big(k_i(t)H_h(t)+h_i(t)H_k(t)\big)\,dt .
\]
In the block with input $j$ at $s$ and output $i$ at $t$, subtract the
kernel $p_0DY_j(s)Y_i(t)$ on the full square. Indeed the Hessian kernel off
the diagonal, with both fields evaluated at $q$, is
\begin{equation}\label{kernel}
 H_{ij}(t,s)=
 \begin{cases}
 p_0DY_i(t)Y_j(s), & s<t,\\
 p_0DY_j(s)Y_i(t), & t<s .
 \end{cases}
\end{equation}
In a spatial basis $(e_\ell)$ the subtracted kernel is
\[
 S_{ij}(t,s)=\sum_{\ell=1}^5
       (p_0DY_j(s)e_\ell)\,Y_i^\ell(t),\qquad \rank S_{ij}\le5 .
\]
The block $H_{ij}$ is a restriction and projection of the finite-rank
operator $B$. Hence
\[
 H_{ij}-S_{ij}={\bf1}_{s<t}\,p_0[Y_j(s),Y_i(t)],\quad
 \rank(H_{ij}-S_{ij})\le\rank B+5 .
\]
Its triangular kernel is continuous off the diagonal and bounded.
Lemma~\ref{causal} gives $\tau_{ij}<1$; take $\tau=\max_{i,j}\tau_{ij}$.
By continuity, taking the later time to $1$ and using antisymmetry gives
\[
 p_0[Y_i(1),Y_j(s)]=0\qquad(\tau<s<1,\ \forall i,j).
\]

Now fix $Z=Y_i(1)$ and apply the linear functional $j^1U\mapsto p_0[Z,U]$
to the two jet identities above. Its values on both $j^1Y_j(t)$ vanish
throughout the final interval. Write $b_Z(t)=p_0[Z,B(t)](q)$.
Differentiating the two zero pairings now yields, almost everywhere,
\[
 0=-u_2(t)b_Z(t),\qquad 0=u_1(t)b_Z(t).
\]
If $b_Z(1)\ne0$, continuity supplies $\sigma<1$ such that
$|b_Z(t)|\ge |b_Z(1)|/2$ for $\sigma<t<1$. Both controls would then
vanish a.e.\ there, contradicting the choice of the last moving piece.
Consequently $b_Z(1)=0$. Transporting this equality to the endpoint for
each choice of $Z$ gives
\[
 p[X_1,[X_1,X_2]](y)=p[X_2,[X_1,X_2]](y)=0 .
\]
Together with the first two layers this proves the assertion.
\end{proof}

\begin{proposition}[The intrinsic Hessian line]\label{line}
Outside a fixed null set, if $m=2$, $y\ne x$, $\rank A=3$, and the minimizer
is unsupported, the image of the intrinsic Hessian spans exactly a line
in $T_yM/W$. Its inverse image is $D_3(y)$, and its nonzero scalar component
has infinite rank on $K$.
\end{proposition}
\begin{proof}
By Lemma~\ref{flags} and \eqref{growth}, $W=D_2(y)$. Write the two
symmetric components of the Hessian as operators $B,C:K\to K^*$.
Set $L=\mathrm{span}\{(B(h,k),C(h,k)):h,k\in K\}$ in quotient coordinates.
Lemma~\ref{scalar} gives $Bh\wedge Ch=0$ for every $h$.
If $\rank B\ge2$, choose $Bh_1,Bh_2$ independent and write
$Ch_i=\lambda_i Bh_i$. Then
\[
 0=B(h_1+h_2)\wedge C(h_1+h_2)
   =(\lambda_2-\lambda_1)Bh_1\wedge Bh_2 .
\]
Write $\lambda=\lambda_1=\lambda_2$. If $Bh\ne0$, choose $i$ with
$Bh\wedge Bh_i\ne0$ and repeat the calculation to get $Ch=\lambda Bh$.
If $Bh=0$, dependence at $h+h_i$ gives
$Ch\in\R Bh_1\cap\R Bh_2=\{0\}$. Thus $C=\lambda B$.
The same applies with $B,C$ interchanged.

In the remaining case each nonzero row has rank one. Symmetry gives
$B=b\,\ell\otimes\ell$, $C=c\,m_0\otimes m_0$ for nonzero rows:
indeed $Bh=d(h)\ell$ implies $d(h)\ell(k)=d(k)\ell(h)$, so $d\propto\ell$.
If both rows are nonzero and $\ell,m_0$ are independent, choose $h$ outside
their two kernels. Then
\[
 Bh\wedge Ch=bc\,\ell(h)m_0(h)\,\ell\wedge m_0\ne0 .
\]
This proves $\dim L\le1$, including a zero row.

Let $\pi:T_yM\to T_yM/W$. If both rows had finite rank,
Lemma~\ref{flags} applied to the two independent quotient coordinate
covectors would give $D_3(y)\subset W$, contrary to
$\dim D_3(y)\ge4>3=\dim W$. Hence $\dim L=1$ and its nonzero
component has infinite rank. For $0\ne p\in L^\perp$ in the quotient dual,
\[
 (p\pi)D^2G|_{K\times K}=0
 \ \Longrightarrow\ D_3(y)\subseteq\ker(p\pi)=\pi^{-1}L,\qquad
 \dim\pi^{-1}L=3+1=4 .
\]
Together with $\dim D_3(y)\ge4$ this gives $D_3(y)=\pi^{-1}L$.
\end{proof}

\proofsection{Uniform two-point correction}\label{repair}
\begin{lemma}[Uniform horizontal correction]\label{two-point}
If a basis of $T_yM$ consists of values of bracket words in a local
orthonormal frame of orders at most $r$, all pairs $z,z'$ in one
smaller coordinate ball can be joined horizontally with length bounded by
$C|z-z'|^{1/r}$. In particular,
\begin{equation}\label{connector}
 d(z,z')\le C |z-z'|^{1/r}
\end{equation}
with $C$ independent of both points. The pointed distance $V$ is finite
and continuous on $M$.
\end{lemma}
\begin{proof}
We give the local construction; see also Jean~\cite{jean}.
Shrink nested coordinate balls so the finitely many flow words below exist
for every base point and all sufficiently small times. In order one,
\[
 F_t(z)=z+tX_t(z),\quad X_t(z)=\int_0^1 X(F_{\theta t}(z))\,d\theta,\quad X_0=X .
\]
For the induction assume $F_s(z)=z+s^kX_s(z)$, $G_t(z)=z+t^lY_t(z)$
with smooth coefficients. Put $H=G_t\circ F_s$, $B=F_s\circ G_t$.
Taking differences in the spatial variable gives $H(z)-B(z)=s^kt^l R(z,s,t)$, where
\[
\begin{aligned}
R={}&\int_0^1 D Y_t(z+\theta s^kX_s(z))X_s(z)\,d\theta\\
 &-\int_0^1 D X_s(z+\theta t^lY_t(z))Y_t(z)\,d\theta .
\end{aligned}
\]
With $B=F_s\circ G_t$ as above, the exact difference is
\[
 B^{-1}H(z)-z
 =s^kt^l\int_0^1 DB^{-1}\big(B(z)+\theta(H(z)-B(z))\big)\,R(z,s,t)\,d\theta .
\]
Shrink the coordinate balls so this segment stays in the domain of $B^{-1}$.
All coefficients are smooth and $DB^{-1}=I$ at $(s,t)=(0,0)$. Thus, on setting
$s=t$,
\[
 \begin{aligned}
 B^{-1}H(z)&=z+t^{k+l}Z(z,t),\\
 Z(z,0)&=DY_0(z)X_0(z)-DX_0(z)Y_0(z)=[X_0,Y_0](z).
 \end{aligned}
\]

Write any resulting word as $W_t(z)=z+t^k Z(z,t)$. Put $t=|s|^{1/k}$. For $s\ne0$ the signed displacement $D$ is
\[
 D(z,s)=
 \begin{cases}
 W_t(z)-z=s Z(z,t),&s>0,\\
 W_t^{-1}(z)-z=s Z(W_t^{-1}(z),t),&s<0 .
 \end{cases}
\]
Writing $D=sZ_\pm$ on the two sides and $D(z,0)=0$,
\[
 \partial_sD=Z_\pm+(t/k)\partial_tZ_\pm\to Z(z,0),\qquad
 D_zD=sD_zZ_\pm\to0 .
\]
The limits are locally uniform; $D(z,s)/s\to Z(z,0)$ gives the derivative
at zero. Thus the signed words are jointly $C^1$ even for even $k$.

Compose five such signed words whose bracket values form a basis at
$y$, of orders at most $r$. Let $F(z,s)$ be their composite endpoint minus $z$, and
$\Psi(z,s)=(z,F(z,s))$. Writing $A=[X_{I_1}(y),\ldots,X_{I_5}(y)]$,
\[
 F(z,0)=0,\quad D\Psi(y,0)=
 \begin{pmatrix}I&0\\0&A\end{pmatrix},\qquad \det A\ne0 .
\]
Write $\Psi^{-1}(z,w)=(z,S(z,w))$. On a smaller product of balls,
\[
 S(z,0)=0,\qquad
 |S(z,w)|\le\sup\|D_wS\|\,|w|,\qquad w=z'-z .
\]
For $s=S(z,z'-z)$ with $|s|\le1$, and $n_i$ generator letters in word
$i$ of order $k_i\le r$, the resulting curve has
\[
 \ell\le\sum_{i=1}^5 n_i |s_i|^{1/k_i}
 \le\Big(\sum_i n_i\Big)|s|^{1/r}
 \le\Big(\sum_i n_i\Big)(\sup\|D_w S\|)^{1/r}|z'-z|^{1/r}.
\]
This gives the two-point estimate.

At each point bracket generation supplies a finite basis. Thus each
horizontal accessibility class is open. Curves can be reversed and
concatenated, so these classes partition $M$, which is connected.
All points are accessible from $x$. The triangle inequality gives
$|V(z)-V(z')|\le d(z,z')$, proving continuity by \eqref{connector}.
\end{proof}

\proofsection{Two missing endpoint directions}\label{two}
\begin{lemma}[Countable straightening charts]\label{straighten}
There is a countable collection of target charts, fixed by the distribution
and metric, with the following property. At a point satisfying
Lemma~\ref{regular-growth}, let $W$ be a three-plane containing $\Delta_y$.
If $m=2$ and $W=D_2(y)$ with $\dim D_3(y)=4$, one chart has a coordinate
axis $e$ equal throughout the chart to a length-three word field with
$e(y)\notin W$, and a second coordinate axis $f(y)\notin D_3(y)$.
If $m=3$, one chart has an axis $e$ equal to a word field of order at
most two with $e(y)\notin W$, and an axis $f(y)\notin W+\R e(y)$.
\end{lemma}
\begin{proof}
Choose a countable frame cover $(U_\nu,X^\nu)_{\nu\in\mathbb N}$. For each
word $|I|\le3$, second countability and the flow-box theorem give
\[
 U_\nu\cap\{X_I^\nu\ne0\}=\bigcup_{j\in\mathbb N}U_{\nu Ij},
 \qquad(\chi_{\nu Ij})_*X_I^\nu=e_1 .
\]
This countable list precedes the choice of $W$. At a point in the statement,
\[
 \begin{array}{lll}
 m=2:&D_2(y)=W\subsetneq D_3(y)
      &\Rightarrow\ \exists |I|=3:\ X_I^\nu(y)\notin W,\\
 m=3:&\dim W=3<\dim D_2(y)
      &\Rightarrow\ \exists |I|\le2:\ X_I^\nu(y)\notin W .
 \end{array}
\]
In the selected chart write $e=e_1$ and identify $W$ with its coordinate
image. Since $\dim(W+\R e)=4$, some coordinate vector $f=e_k$ lies outside
$W+\R e$. Thus $T_yM=W\oplus\R e\oplus\R f$ in these coordinates.
\end{proof}

\begin{lemma}[Quadratic and cubic alternative]\label{polynomials}
Let $Q,P$ be homogeneous polynomials of degrees two and three on a real
vector space $K$, given by symmetric multilinear forms. Suppose $Q$ has
rank at least three. Either $dQ(h),dP(h)$ are independent somewhere, or
$P$ vanishes identically.
\end{lemma}
\begin{proof}
In the dependent case apply $dQ\wedge dP=0$ to the radial vector $h$
and an arbitrary $k$. Euler's identities give
\[
 0=(dQ\wedge dP)_h(h,k)=2Q(h)dP_h(k)-3P(h)dQ_h(k).
\]
Include any prescribed vector of $K$ in a finite-dimensional subspace
on which $Q$ has rank at least three. Such a subspace exists by adjoining
the vectors of a nonsingular minor of the bilinear form. On a ball in
that subspace with $Q\ne0$,
\[
 d(P^2/Q^3)=\frac{P}{Q^4}\bigl(2Q\,dP-3P\,dQ\bigr)=0,
\]
so
$P^2=bQ^3$ on the ball and by polynomial identity on the whole subspace.
For linear forms $\ell_1,\ell_2$,
\[
 Q=\ell_1\ell_2\ \Longrightarrow\
 D^2Q=\ell_1\otimes\ell_2+\ell_2\otimes\ell_1
 \ \Longrightarrow\ \rank Q\le2 .
\]
Homogeneity therefore makes $Q$ irreducible on this finite slice.
In its real polynomial ring, unique factorization gives
\[
 b\ne0\ \Longrightarrow\
 2\operatorname{ord}_Q(P)=\operatorname{ord}_Q(bQ^3)=3 .
\]
Thus $b=0$, so $P=0$ on the slice, including at the prescribed vector.
\end{proof}

\begin{proposition}[A common null set in endpoint rank three]\label{two-missing}
Fix $x$. For $m=2$ or $3$ there is a Borel null set $N_{x,3}$
such that, for every $y\notin N_{x,3}$ distinct from $x$, every minimizing
control from $x$ to $y$ with endpoint derivative of rank three has a
supporting length covector.
\end{proposition}
\begin{proof}
Let $N_{\rm gr},N_{\rm scalar}$ be Borel null covers from
Lemmas~\ref{regular-growth} and~\ref{scalar}. Use the countable index set
\[
 \begin{aligned}
 \Lambda&=\{(\chi,e,f):\chi\text{ in Lemma~\ref{straighten}},\
                  e\ne f\text{ coordinate axes}\}\\
         &\hspace{12mm}\times\{(2,2),(2,3),(2,4)\}.
 \end{aligned}
\]
For each $\lambda\in\Lambda$ take the corresponding $N_\lambda$ in
Lemma~\ref{planes}. Thus
\[
 N_{x,3}=N_{\rm gr}\cup N_{\rm scalar}\cup\bigcup_{\lambda\in\Lambda}N_\lambda,
 \qquad \mu(N_{x,3})\le0+0+\sum_{\lambda\in\Lambda}0=0 .
\]
Fix $y\notin N_{x,3}$, $y\ne x$ and an unsupported minimizer with $\rank A=3$.
Lemma~\ref{kernel-slice} transports the support condition and Hessian to
the selected chart.

For $m=2$, Proposition~\ref{line} and Lemma~\ref{straighten} give
\[
 W=D_2(y),\quad D_3(y)=W\oplus\R e,\quad T_yM=W\oplus\R e\oplus\R f .
\]
On the slice write $G(S(k))-y=F_1(k)e+F_2(k)f$. Then
\[
 F_i(0)=0,\quad DF_i(0)=0,\quad \rank D^2F_1(0)=\infty,\quad D^2F_2(0)=0 .
\]
Taylor expansion on finite-dimensional direction families gives, locally
uniformly in $h$,
\begin{equation}\label{jet-scales}
 \begin{aligned}
 F_1(th)&=t^2Q(h)+O(t^3), & Q(h)&=\tfrac12 D^2F_1(0)[h,h],\\
 F_2(th)&=t^3P(h)+O(t^4), & P(h)&=\tfrac16 D^3F_2(0)[h,h,h].
 \end{aligned}
\end{equation}

If the rows are independent at $h_*$, choose $b_1,b_2$ with
\[
 M(v)=\det
 \begin{pmatrix}dQ_v b_1&dQ_v b_2\\ dP_v b_1&dP_v b_2\end{pmatrix},
 \qquad M(h_*)\ne0 .
\]
For any strictly shorter direction $h_0$ the polynomial
$M(h_0+\lambda h_*)$ has leading coefficient $M(h_*)$ at degree three.
Hence some arbitrarily small $\lambda$ gives $M(h_0+\lambda h_*)\ne0$.
Take such $h$ in the open shorter set. Lemmas~\ref{scaled-box}
and~\ref{support} give one open rectangle $U$ and $c,t_0>0$ with
\[
 \begin{aligned}
 &\forall (a,b)\in U,\ 0<t<t_0\quad\exists h'\text{ near }h:\\
 &(F_1(th'),F_2(th'))=(t^2a,t^3b),\\
 &V(y+t^2ae+t^3bf)\le J(S(th'))\le V(y)-ct .
 \end{aligned}
\]
This is precisely the $(2,3)$ test of Lemma~\ref{planes}.

Otherwise $P=0$ everywhere by Lemma~\ref{polynomials}, since
Proposition~\ref{line} gives infinite rank for $Q$.

Choose a strictly shorter $h$ with $dQ(h)b_1\ne0$ for some $b_1$;
the open shorter set cannot lie in the radical of $Q$. For $a$ near
$Q(h)$ the scaled implicit function theorem solves
$F_1(t(h+\lambda b_1))=t^2a$ smoothly for $\lambda(t,a)$.
Put $v_a=h+\lambda(0,a)b_1$ and $C(a)=D^4F_2(0)[v_a,v_a,v_a,v_a]/24$.
To obtain a remainder uniform in $a$, take a compact interval $I$
about $Q(h)$ inside the implicit-function neighborhood and put
$v(t,a)=h+\lambda(t,a)b_1$. The derivatives of $F_2$ through order
three vanish at zero on this slice. Integral Taylor expansion gives
\[
 \frac{F_2(tv(t,a))}{t^4}
 =\frac1{3!}\int_0^1(1-s)^3
 D^4F_2(stv(t,a))[v(t,a)^4]\,ds ,
\]
Here $[v^4]$ denotes four copies of $v$. Uniform continuity on
$[0,1]\times[0,t_0]\times I$ gives the estimate
\[
 \begin{aligned}
 &\sup_{a\in I}|t^{-4}F_2(tv(t,a))-C(a)|\\
 &\quad\le\tfrac1{24}\sup_{\substack{a\in I\\0\le s\le1}}
 \big|D^4F_2(stv(t,a))[v(t,a)^4]-D^4F_2(0)[v_a^4]\big|
 \longrightarrow0 .
 \end{aligned}
\]
Here $(3!)^{-1}\int_0^1(1-s)^3\,ds=1/24$. Thus
\begin{equation}\label{fourth-output}
 F_2(t(h+\lambda(t,a)b_1))=t^4(C(a)+r_t(a)),\qquad
 \sup_{a\in I}|r_t(a)|\longrightarrow 0 .
\end{equation}
Let $z_t$ be this endpoint, let $z'_t=y+t^2ae+t^4bf$, and denote
the constant in \eqref{connector} for $r=4$ by $K_0$. The coordinate
vector $f$ has unit Euclidean norm. The two bounds on cost are
\begin{align}
 V(z_t)&\le V(y)-c_0t, \label{fourth-saving}\\
 d(z_t,z'_t)&\le K_0\,t\,|b-C(a)-r_t(a)|^{1/4}. \label{fourth-cost}
\end{align}
Choose $\rho>0$ with $K_0(3\rho)^{1/4}\le c_0/2$. Shrink $I$ until
$|C(a)-C(Q(h))|<\rho$ throughout. Choose $|b-C(Q(h))|<\rho$ and then
$t$ small enough for $\sup_I|r_t|<\rho$. Adding
\eqref{fourth-saving} and \eqref{fourth-cost} gives, throughout the rectangle,
\begin{equation}\label{box24}
 V(y+t^2ae+t^4bf)\le V(z_t)+d(z_t,z'_t)\le V(y)-\tfrac12 c_0t .
\end{equation}
This calculation also applies when $C(Q(h))=0$.

In both alternatives the constructed rectangle places $y$ in the set of
Lemma~\ref{planes}, contrary to $y\notin N_{x,3}$. The axes are those of
Lemma~\ref{straighten} even if the minimizing control was written using
different frames. Only the rectangle and its constants depend on that control.

For $m=3$, $\rank A=3$, use Lemmas~\ref{regular-growth} and~\ref{straighten}:
\[
 \Delta_y=W,\quad \dim D_2(y)\ge4,\quad D_3(y)=T_yM,\quad
 T_yM=W\oplus\R e\oplus\R f .
\]
Here $e$ is a straightened word of order at most two outside $W$.
For order two apply the signed word from Lemma~\ref{two-point} with
time $\sqrt{|s|}$, using the inverse word for $s<0$.
If $w_s$ is the endpoint of this word from $z$ and $V_0$ is a
smaller coordinate ball, smoothness of the word on its buffered chart
gives constants $A_0,B_0$ uniform for $z\in V_0$. Using
\eqref{connector} with $r=3$ gives
\[
 \begin{aligned}
 |w_s-(z+se)|&\le B_0|s|^{3/2},\qquad d(z,w_s)\le A_0|s|^{1/2},\\
 d(z,z+se)&\le A_0|s|^{1/2}+K_0'(B_0|s|^{3/2})^{1/3}\\
           &=(A_0+K_0'B_0^{1/3})|s|^{1/2}.
 \end{aligned}
\]
For an order-one
field its flow gives the exact translation at cost $|s|$.

A shorter slice ray ends at $y+F(t)e+H(t)f$ with length at most $V(y)-ct$.
Both functions start with zero value and derivative. If $H''(0)=0$, the
quotients $F(t)/t^2,H(t)/t^3$ converge. For $(a,b)$ in a sufficiently
small box around their limits, move to $y+t^2ae+t^3bf$ by first using
the pure $e$ displacement and then \eqref{connector} on the remainder.
With constants $C_2,C_3$ uniform in the starting point, the cost divided
by $t$ is bounded by
\[
 C_2|a-F(t)/t^2|^{1/2}+C_3|b-H(t)/t^3|^{1/3}.
\]
Choose a box of radius $\eta$ about the limits with
$C_2(2\eta)^{1/2}+C_3(2\eta)^{1/3}<c/2$, then take $t$ small enough
that both quotients are within $\eta$ of their limits.
The length bound after both corrections is therefore $V(y)-ct/2$.

If $H''(0)\ne0$, write $\alpha=F''(0)/2$, $\beta=H''(0)/2$ and put
$s=t\theta$. Taylor expansion uniformly for $1/2\le\theta\le3/2$ gives
\[
 t^{-2}F(t\theta)=\alpha\theta^2+o(1),\qquad
 t^{-2}H(t\theta)=\beta\theta^2+o(1).
\]
First choose $\delta>0$ with $C_2\sqrt{3\delta}<c/4$, then choose
$0<\epsilon<1/2$ with $|\alpha| (2\epsilon+\epsilon^2)<\delta$.
Put $\theta_\pm=1\pm\epsilon$ and
\[
 \rho=\tfrac12\min_{\pm}|\beta\theta_\pm^2-\beta|>0 .
\]
Choose $t$ small enough that the two uniform errors above are less than
$\delta$ and $\rho$, respectively. Since $\beta\theta_-^2$ and
$\beta\theta_+^2$ lie on opposite sides of $\beta$, continuity gives
\[
 \forall b\in(\beta-\rho,\beta+\rho)\ \exists\theta\in[\theta_-,\theta_+]:
 \quad H(t\theta)=t^2b .
\]
For every such $\theta$ and $|a-\alpha|<\delta$ we have
\[
 |a-t^{-2}F(t\theta)|
 \le |a-\alpha|+|\alpha(1-\theta^2)|
      +|\alpha\theta^2-t^{-2}F(t\theta)|<3\delta .
\]
Thus for $z=y+F(t\theta)e+H(t\theta)f$ and $z'=y+t^2ae+t^2bf$,
\[
 \begin{aligned}
 d(z,z')&\le C_2t|a-t^{-2}F(t\theta)|^{1/2}<ct/4,\\
 V(z')&\le V(z)+d(z,z')\le V(y)-ct\theta+ct/4
          \le V(y)-ct/4 .
 \end{aligned}
\]
The intervals for $a,b$ were fixed before $t$, giving the $(2,2)$ box.
Lemma~\ref{planes} again puts $y$ in $N_{x,3}$, a contradiction.
This exhausts the two horizontal ranks in the proposition.
\end{proof}

\proofsection{Critical normal parameters and completion of the proof}
\begin{lemma}[Supported singular paths give critical normal parameters]\label{normal}
Endpoints of supported singular minimizers from $x$ form a set of zero
smooth outer volume.
\end{lemma}
\begin{proof}
Write the consecutive normalized pieces on $[0,T]$. Transport a supporting
row $\eta$ backwards by $\dot p=-\sum_i u_i p DX_i$, taking the cotangent
transition at each joint. Variation of constants in the state equation
gives $\eta Av=\int_0^T a\cdot v$ with $a_i(t)=p(t)X_i(q(t))$.
With $\nu=|u|$, the two tests $v=0,2u$ give
\[
 0\ge J(u)-\eta Au,\qquad 2J(u)\ge J(u)+\eta Au
 \quad\Longrightarrow\quad \eta Au=J(u).
\]
Consequently $\int a\cdot v\le\int|v|$ for all $v$. Put
$B=\{|a|>1\}$ and take $v={\bf1}_B\,a/|a|$ (zero off $B$). Then
$\int_B(|a|-1)\le0$, so $|B|=0$.
The support equality and the pointwise norm bound now give
\begin{equation}\label{deficit}
 \int_0^T (|u|-a\cdot u)\,dt=J(u)-\eta Au=0,\qquad
 |u|-a\cdot u\ge0\quad\text{a.e.}
\end{equation}
Where $u\ne0$, equality forces $a=u/|u|$ by Cauchy--Schwarz. Hence
$u=\nu a$ almost everywhere for $\nu=|u|$, including at pauses. Consequently
\[
  (\dot q,\dot p)=\nu(t)\vec H(q,p),\qquad
  H(q,p)=\tfrac12\sum_i\big(pX_i(q)\big)^2 .
\]
In a frame the Hamilton equations read
\[
 \partial_p H=\sum_i a_iX_i,\qquad
 -\partial_q H=-\sum_i a_i p DX_i,\qquad
 \sum_i a_i^2=\|a\|^2 .
\]
The last expression is invariant under orthogonal frame changes. For
$L=\int_0^T\nu=0$, $q(t)=x$. For $L>0$ put
$\sigma(t)=L^{-1}\int_0^t\nu$, $\xi_0=Lp(0)$, and
$S_L(q,p)=(q,Lp)$. Quadratic homogeneity gives
\[
 \vec H(S_L(q,p))=L\,DS_L\,\vec H(q,p).
\]
Thus the following equality holds by uniqueness of the integral equation
on successive coordinate intervals, also on intervals where $\sigma'=0$:
\[
 S_L(q(t),p(t))=\Phi_H^{\sigma(t)}(x,\xi_0),\qquad
 \sigma(0)=0,\quad \sigma(T)=1 .
\]
Here the orbit extends along the range $[0,1]$ of the clock: up to any
putative first finite escape time the orbit agrees with the compact
reference curve; the local continuation theorem then extends it.
Write $P(\xi)=\pi_M\Phi_H^1(x,\xi)$ on its open domain in $T_x^*M$.
For nearby $\xi$, let $(q^\xi,p^\xi)$ be the orbit and keep the same clock.
In each of the original charts use controls
\[
   v_i(\xi,t)=\frac{\nu(t)}{L}p^\xi(\sigma(t))X_i(q^\xi(\sigma(t))).
\]
Let $b_i(\xi,s)=p^\xi(s)X_i(q^\xi(s))$ on the compact clock interval of
each piece. On a smaller parameter ball its first two parameter derivatives
are bounded. Taylor's formula gives, with a uniform constant over the pieces,
\[
 \begin{aligned}
 Dv_i(\xi)[h](t)&=\sigma'(t)D_\xi b_i(\xi,\sigma(t))[h],\\
 \|v(\xi+h)-v(\xi)-Dv(\xi)[h]\|_{\mathcal U}
       &\le C\|\sigma'\|_2 |h|^2,\qquad v(\xi_0)=u .
 \end{aligned}
\]
By uniqueness on each consecutive piece the nodes match. In terminal
coordinates this gives
\[
 P=G\circ v,\qquad DP(\xi_0)=DG(u)Dv(\xi_0),\qquad
 \rank DP(\xi_0)\le\rank DG(u)<5 .
\]
On the open time-one domain in $T_x^*M$, classical Sard gives
\[
 \mu^*(P(\operatorname{Crit}P))=0,\qquad \mu^*(\{x\})=0 .
\]
These two sets cover all supported singular minimizing endpoints at $x$.
\end{proof}

\begin{proof}[Proof of Theorem~\ref{main}]
Fix $x$ and let $\Sigma_x$ denote the endpoint set in the statement.
Choose Borel null covers $N_{\rm gr}$ of the growth exception in
Lemma~\ref{regular-growth} and $N_{x,4}$ of the corank-one exception
in Lemma~\ref{scalar}, using a countable measuring atlas. Take
$N_{x,3}$ from Proposition~\ref{two-missing} when $m=2,3$ and the
empty set when $m=4$. Let $P_x$ be the normal time-one endpoint map
of Lemma~\ref{normal} and let $\operatorname{Crit}P_x$ denote its critical
parameters in its open domain. We claim the set inclusion
\begin{equation}\label{final-inclusion}
 \Sigma_x\subset \{x\}\cup N_{\rm gr}\cup N_{x,4}\cup N_{x,3}
                  \cup P_x(\operatorname{Crit}P_x).
\end{equation}
Take a witness $u$ for $y\in\Sigma_x$ outside the first four sets.
Write $r=\rank DG(u)<5$. For $m=1$, $[aX,bX]=(aXb-bXa)X$ contradicts
bracket generation. For $m=2,3,4$, Lemmas~\ref{flags} and~\ref{regular-growth} give
\[
 \begin{array}{c|c|c}
 m & \text{lower bound for }r &\text{remaining }r\\ \hline
 2 & r\ge\dim D_2(y)=3 &3,4\\
 3 & r\ge\dim\Delta_y=3&3,4\\
 4 & r\ge\dim\Delta_y=4&4 .
 \end{array}
\]
By Lemma~\ref{scalar} when $r=4$ and Proposition~\ref{two-missing} when $r=3$,
\[
 \left.
 \begin{array}{l}
 r=4,\ y\notin N_{x,4}\\[-2pt]
 \text{or }r=3,\ y\notin N_{x,3}
 \end{array}\right\}
 \Longrightarrow\mathcal S(u)
 \overset{\text{Lemma~\ref{normal} (proof)},\,y\ne x}{\Longrightarrow}
 y\in P_x(\operatorname{Crit}P_x).
\]
This proves \eqref{final-inclusion}. Smooth coordinate densities and a
countable chart cover carry the coordinate Sard null set to a $\mu$-null set.
By outer subadditivity,
\[
 \mu^*(\Sigma_x)\le\mu^*(\{x\})+\mu(N_{\rm gr})+\mu(N_{x,4})
          +\mu(N_{x,3})+\mu^*(P_x(\operatorname{Crit}P_x))=0 .
\]
\end{proof}

\begin{thebibliography}{9}
\bibitem{strong}
A. Belotto da Silva, A. Figalli, A. Parusi\'nski and L. Rifford,
Strong Sard conjecture and regularity of singular minimizing geodesics
for analytic sub-Riemannian structures in dimension 3,
\emph{Invent. Math.} \textbf{229} (2022), 395--448.
\href{https://doi.org/10.1007/s00222-022-01111-2}{doi:10.1007/s00222-022-01111-2}.
\bibitem{huang}
Y. Huang, Autonomous disproofs of the sum-product conjecture over
$\mathbb R$ with GPT-5.5 Pro, \href{https://arxiv.org/abs/2607.20525}{arXiv:2607.20525} (2026).
\bibitem{jean}
F. Jean, Uniform estimation of sub-Riemannian balls,
\emph{J. Dynam. Control Systems} \textbf{7} (2001), 473--500.
\href{https://doi.org/10.1023/A:1013154500463}{doi:10.1023/A:1013154500463}.
\bibitem{rifford}
L. Rifford, Subdifferentials and minimizing Sard conjecture in sub-Riemannian
geometry, \emph{J. \'{E}c. polytech. Math.} \textbf{10} (2023), 1195--1244.
\href{https://doi.org/10.5802/jep.242}{doi:10.5802/jep.242}.
\end{thebibliography}
\end{document}
